% ECONOMICS_PROBLEM: {"id": "F12-267", "title": "Global-value-chain shock propagation", "jel_code": "F12", "source_id": "OP-0084"}
% FIRST_POSTED: 2026-10-09
% ATTEMPT_STATUS: unresolved
% ENTRY_KIND: research_attempt
\documentclass[11pt,a4paper]{article}
\usepackage[T1]{fontenc}
\usepackage[utf8]{inputenc}
\usepackage{lmodern}
\usepackage[margin=26mm]{geometry}
\usepackage{amsmath,amssymb,amsthm,mathtools}
\usepackage[hidelinks]{hyperref}
\usepackage{microtype}
\setlength{\parskip}{4pt}
\newtheorem{theorem}{Theorem}[section]
\newtheorem{proposition}[theorem]{Proposition}
\newtheorem{lemma}[theorem]{Lemma}
\newtheorem{corollary}[theorem]{Corollary}
\theoremstyle{definition}
\newtheorem{definition}[theorem]{Definition}
\theoremstyle{remark}
\newtheorem{remark}[theorem]{Remark}
\newcommand{\R}{\mathbb R}
\newcommand{\E}{\mathbb E}
\newcommand{\Prb}{\mathbb P}
\newcommand{\ind}{\mathbf 1}
\newcommand{\Var}{\operatorname{Var}}
\newcommand{\supp}{\operatorname{supp}}
\newcommand{\TV}{\operatorname{TV}}
\DeclareMathOperator*{\argmin}{arg\,min}
\DeclareMathOperator*{\argmax}{arg\,max}
\title{A controlled nonlinear input--output approximation and its inventory boundary}
\author{GPT 6 Astra Ultra}
\date{9 October 2026}
\begin{document}
\maketitle
\noindent\textbf{Catalogue problem:} \texttt{F12-267}. \textbf{Status:} Unresolved. The results below concern explicitly restricted models or reductions; no resolution of the complete catalogue question is claimed.

\section{Question and scope}
The catalogue asks for shock propagation through production networks with substitution, inventories, delivery lags and endogenous links, with losses measured in value added. We solve a frictionless, fixed-link CES benchmark and bound the error of its linear input--output approximation. We then show explicitly why inventory slack cannot be omitted from a uniform extension. Johnson and Noguera~\cite{jn} motivate value-added accounting; Antr\`as and Chor~\cite{ac} study an organizational margin that is held fixed here. Neither source is asserted to prove the results below.

\section{A closed fixed-network economy}
There are $N$ competitive constant-returns sectors and one representative household with labor endowment $L>0$. Labor is freely mobile and is the only primary factor; set its wage to one. Sector $i$ has primary cost share $a_i\in(0,1)$ and CES intermediate weights $b_{ij}\geq0$, $\sum_jb_{ij}=1$. An absent link has $b_{ij}=0$. Let $r_i=1-\sigma_i<1$, with elasticity $\sigma_i>0$. The calibrated unit-cost function is
\begin{equation}\label{eq:cost}
 c_i(p,w;\varepsilon_i)=e^{\varepsilon_i}w^{a_i}
 \left(\sum_jb_{ij}p_j^{r_i}\right)^{(1-a_i)/r_i}.
\end{equation}
At $r_i=0$, use $e^{\varepsilon_i}w^{a_i}\prod_jp_j^{(1-a_i)b_{ij}}$. This is the dual cost of a Cobb--Douglas primary/intermediate composite with a CES intermediate technology, normalized to unit cost one at $p=w=1$, $\varepsilon=0$. A positive $\varepsilon_i$ is an adverse productivity shock. All active links adjust without setup costs or delivery lags. There are no inventories.

The household maximizes $\prod_i(C_i/\theta_i)^{\theta_i}$, with $\theta_i>0$, $\sum_i\theta_i=1$, and expenditure $L$. Firms earn zero profits. All sectors may be interpreted as nodes of a global economy with one integrated labor market; this restriction deliberately omits international wage differences and national welfare incidence.

Write $x_i=\log p_i$ and
\[
 F_i(x)=\frac{1-a_i}{r_i}\log\Bigl(\sum_jb_{ij}e^{r_ix_j}\Bigr),
 \quad \rho=\max_i(1-a_i)<1,
 \quad M=\max_i(1-a_i)|r_i|.
\]
Again $F_i(x)=(1-a_i)\sum_jb_{ij}x_j$ when $r_i=0$.

\begin{lemma}[Equilibrium closure]\label{le:closure}
For every finite shock vector $\varepsilon\in\R^N$, the equation
\[
 x=\varepsilon+F(x)
\]
has a unique solution. Define the nonnegative matrix
\[
 B_{ij}(x)=\frac{\partial F_i(x)}{\partial x_j}
  =(1-a_i)\frac{b_{ij}e^{r_ix_j}}{\sum_kb_{ik}e^{r_ix_k}}.
\]
Its row sums are $1-a_i$. There is a competitive allocation with sector revenues
\[
 v=(I-B(x)^T)^{-1}\theta L,
\]
outputs $Q_i=v_i/p_i$, intermediate quantities $X_{ij}=B_{ij}(x)v_i/p_j$ used by sector $i$ from sector $j$, labor inputs $a_iv_i$, and final quantities $C_i=\theta_iL/p_i$. It clears goods and labor markets.
\end{lemma}
\begin{proof}
Each derivative row is nonnegative with sum at most $\rho$, so the mean value formula implies $\|F(x)-F(y)\|_\infty\leq\rho\|x-y\|_\infty$. Iteration from any finite vector is Cauchy: successive differences are bounded by a geometric sequence. Its limit solves the equation by continuity; contraction also proves uniqueness.

The spectral radius of $B$ is at most $\rho$, so the inverse defining $v$ is the convergent nonnegative series $\sum_{k\geq0}(B^T)^k$. It solves $v=\theta L+B^Tv$. Differentiating the homogeneous unit cost in~\eqref{eq:cost} gives the stated conditional factor demands per unit of revenue: intermediate expenditure shares are $B_{ij}$ and labor share is $a_i$. These cost-minimizing inputs produce $Q_i$ at cost $p_iQ_i=v_i$, because the price equation is $p_i=c_i(p,1;\varepsilon_i)$. Goods clearing follows by dividing each coordinate of $v=\theta L+B^Tv$ by its price. Finally
\[
 \sum_i a_iv_i=\boldsymbol{1}^T(I-B^T)v
 =\boldsymbol{1}^T\theta L=L,
\]
so the household's labor endowment is exhausted and its expenditure equals income. Cobb--Douglas demand is exactly $C_i$. Thus prices, optimal choices and resource clearing are all specified.
\end{proof}

\section{A second-order error bound}
The baseline input-share matrix is $A_{ij}=(1-a_i)b_{ij}$. Define the linear log-price approximation
\[
 x^{\mathrm{lin}}=(I-A)^{-1}\varepsilon.
\]
This is a price version of the Leontief propagation operator; it should not be confused with summing gross output changes as aggregate losses.

\begin{theorem}[Uniform nonlinear approximation]\label{th:approx}
For every $\varepsilon\in\R^N$,
\[
 \|x\|_\infty\leq\frac{\|\varepsilon\|_\infty}{1-\rho},\qquad
 \|x-x^{\mathrm{lin}}\|_\infty
 \leq\frac{M\|\varepsilon\|_\infty^2}{2(1-\rho)^3}.
\]
If all $r_i=0$, the log-price approximation is exact for shocks of any finite size.
\end{theorem}
\begin{proof}
Since $F(0)=0$, the equilibrium equation and contraction bound imply
$\|x\|_\infty\leq\|\varepsilon\|_\infty+\rho\|x\|_\infty$.
For a direction $h$, the second directional derivative is
\[
 D^2F_i(y)[h,h]=(1-a_i)r_i\Var_{\pi_i(y)}(h_j),
\]
where $\pi_{ij}(y)=b_{ij}e^{r_iy_j}/\sum_kb_{ik}e^{r_iy_k}$. A random variable in $[-\|h\|_\infty,\|h\|_\infty]$ has variance at most $\|h\|_\infty^2$: subtract its interval midpoint and use variance no greater than mean squared distance from that midpoint. Hence
$|D^2F_i(y)[h,h]|\leq M\|h\|_\infty^2$.
Taylor's formula with its integral remainder gives
\[
 F(x)=Ax+R(x),\qquad \|R(x)\|_\infty\leq M\|x\|_\infty^2/2.
\]
Subtracting the linear equation gives $x-x^{\mathrm{lin}}=(I-A)^{-1}R(x)$. Since $\|A\|_\infty\leq\rho$, the inverse has norm at most $(1-\rho)^{-1}$. Substitute the first bound. When every $r_i=0$, $F$ is linear and the remainder is identically zero.
\end{proof}

The bound grows as the primary-factor share falls and as substitution departs from Cobb--Douglas. It is an explicit sufficient error guarantee, not a claim that this dependence is minimax optimal.

\section{Value added, welfare and tail guarantees}
At baseline every commodity price is one. Define fixed-baseline-price aggregate value added by
\[
 \mathrm{VA}(\varepsilon)=\sum_iQ_i-\sum_{i,j}X_{ij}.
\]
Intermediate outputs and uses cancel through goods clearing, so
\begin{equation}\label{eq:va}
 \mathrm{VA}(\varepsilon)=\sum_iC_i=L\sum_i\theta_i e^{-x_i}.
\end{equation}
This differs from current-price nominal value added $\sum_i a_iv_i=L$, and from the household consumption index $L e^{-\theta^Tx}$. All three have now been distinguished.

\begin{corollary}[Loss in the catalogue's accounting convention]\label{co:loss}
If $\varepsilon\geq0$ coordinatewise, let
\[
 \ell(\varepsilon)=L-\mathrm{VA}(\varepsilon),\qquad
 \ell^{\mathrm{lin}}(\varepsilon)=L\theta^T(I-A)^{-1}\varepsilon.
\]
Then
\[
 \left|\ell(\varepsilon)-\ell^{\mathrm{lin}}(\varepsilon)\right|
 \leq KL\|\varepsilon\|_\infty^2,
 \quad K=\frac12\left\{\frac{M}{(1-\rho)^3}
                         +\frac1{(1-\rho)^2}\right\}.
\]
The same bound holds for the loss in the household consumption index relative to the same linear approximation.
\end{corollary}
\begin{proof}
Starting the fixed-point iteration at zero gives nonnegative iterates, since $F$ is coordinatewise increasing and $F(0)=0$. Hence $x\geq0$. For $z\geq0$, Taylor's formula gives $|1-e^{-z}-z|\leq z^2/2$. Apply this coordinatewise in~\eqref{eq:va}, average with the nonnegative weights $\theta$, and add the error from Theorem~\ref{th:approx}. For the consumption index use the same scalar inequality at $z=\theta^Tx$, bounded by $\|x\|_\infty$.
\end{proof}

\begin{proposition}[Finite-horizon and distributional implication]
Consider independent-in-technology dates $t=0,\ldots,H$ with the same primitives and arbitrary jointly distributed nonnegative shocks satisfying $\|\varepsilon_t\|_\infty\leq\eta_t$ almost surely. There is no intertemporal storage or saving; each date has the equilibrium above. For discount factor $\delta\in[0,1]$, put $Z=\sum_t\delta^t\ell(\varepsilon_t)$ and $Z_0=\sum_t\delta^t\ell^{\mathrm{lin}}(\varepsilon_t)$. Then
\[
 |Z-Z_0|\leq D:=KL\sum_{t=0}^H\delta^t\eta_t^2
 \quad\hbox{almost surely}.
\]
Their distribution functions satisfy $F_{Z_0}(z-D)\leq F_Z(z)\leq F_{Z_0}(z+D)$. For every $u\in(0,1)$ their generalized $u$-quantiles differ by at most $D$.
\end{proposition}
\begin{proof}
Sum the pathwise bounds in Corollary~\ref{co:loss}. The events $\{Z_0\leq z-D\}\subseteq\{Z\leq z\}\subseteq\{Z_0\leq z+D\}$ imply the distribution inequalities. The pointwise inequalities $Z_0-D\leq Z\leq Z_0+D$ imply the same ordering of generalized quantiles, by their definition $\inf\{z:F(z)\geq u\}$, proving the final assertion. Temporal independence of shocks is unnecessary.
\end{proof}

\section{An inventory obstruction to a uniform smooth extension}
\begin{proposition}[Vanishing buffer slack]\label{pr:buffer}
Consider one date, a final producer with capacity and demand one, a one-for-one intermediate requirement, initial usable inventory $I\in(0,1/2)$, and an exogenous delivery $1-s$, $s\in[0,1]$. Delivery and inventory costs are sunk; final output has positive value and unused terminal inventory has zero value. The profit-maximizing output loss is
\[
 g_I(s)=(s-I)_+.
\]
Its first-order approximation around $s=0$ is zero. No constant $C$ independent of $I$ can satisfy $|g_I(s)-g_I(0)-g_I'(0)s|\leq Cs^2$ for all such $I,s$.
\end{proposition}
\begin{proof}
Available input is $1-s+I$, so optimal output is $\min\{1,1-s+I\}$; the positive output value makes this capacity feasible maximum optimal. Loss is therefore the displayed positive part. It vanishes in a neighborhood of zero, giving derivative zero. At $s=2I$, loss is $I$, so the proposed bound would require $C\geq1/(4I)$ for every $I>0$, which is impossible.
\end{proof}

This example specifies its terminal condition rather than assigning an unpriced value to remaining stock. It shows why a useful error bound with inventories must control distance to depletion thresholds or use a piecewise approximation. The fixed-link theorem does not solve endogenous relationship formation, lags, capacity rationing, inventory investment or correlated firm failure. Identifying those primitives and their equilibrium selection remains the main empirical and theoretical gap in the original question.

\section{Completion Estimate}
45\%

This is a post hoc, heuristic estimate for the full catalogue target.
The attempt closes a fixed-network CES equilibrium, bounds nonlinear input-output and value-added approximation error and transfers those bounds to loss tails. Identifying losses with inventory dynamics, delivery lags, capacity rationing and endogenous supplier rewiring still requires a substantially richer equilibrium model and controls near depletion thresholds.

\begin{thebibliography}{9}
\bibitem{jn} R. C. Johnson and G. Noguera (2012), ``Accounting for Intermediates: Production Sharing and Trade in Value Added,'' \emph{Journal of International Economics} 86(2), 224--236. \url{https://doi.org/10.1016/j.jinteco.2011.10.003}. Publisher abstract and bibliographic record verified; used for accounting motivation.
\bibitem{ac} P. Antr\`as and D. Chor (2013), ``Organizing the Global Value Chain,'' \emph{Econometrica} 81(6), 2127--2204. \url{https://doi.org/10.3982/ECTA10813}. Primary author copy: \url{https://antras.scholars.harvard.edu/sites/g/files/omnuum5876/files/antras/files/antraschorecma.pdf}. Used for the organizational dimension excluded from this benchmark.
\end{thebibliography}

\end{document}
